\section{Technical Appendix: Proofs and Boundary Cases}
\label{app:proofs}

This appendix states the exact scope of every named mathematical result used
by the paper.  Conditional probabilities are asserted only when their
conditioning events have positive probability.  All one-sided results assume
that only an adverse initial decision can change and that a non-appellant
retains the initial decision.

\subsection{Operators and confusion-rate transformations}

\paragraph{Theorem 1 (linear operator preservation criterion).}
Let $V$ be a finite-dimensional vector space of feasible signed
perturbations, let $T_K:V\to V$ be the linear map induced by a fixed
collection of stochastic kernels, and let $L:V\to W$ encode a homogeneous
linear fairness constraint.  The following are equivalent:
\begin{enumerate}
  \item $Lv=0$ implies $LT_Kv=0$ for every $v\in V$;
  \item $\ker L\subseteq\ker(LT_K)$;
  \item there is a linear map $M:\operatorname{im}L\to W$ such that
        $LT_K=ML$ on $V$.
\end{enumerate}

\begin{proof}
The first two conditions are identical.  Under the second, define
$M(Lv)=LT_Kv$.  If $Lv=Lw$, then $v-w\in\ker L$, so
$LT_K(v-w)=0$ and the definition does not depend on the representative.
It is linear by construction.  Conversely, $LT_K=ML$ immediately implies
$LT_Kv=0$ whenever $Lv=0$.
\end{proof}

For probability models, $V$ must be the span of feasible signed differences.
An assertion made only on a thin or boundary-only subset of probability
vectors does not automatically extend to the full algebraic kernel.  Affine
constraints must be homogenized with a constant coordinate.

\paragraph{Proposition 1 (one-sided confusion-rate transformation).}
For every group $g$ with defined label-conditional rates,
\begin{align*}
 \mathrm{TPR}'_g&=\mathrm{TPR}_g+\mathrm{FNR}_g\kappa_g^+,&
 \mathrm{FNR}'_g&=\mathrm{FNR}_g(1-\kappa_g^+),\\
 \mathrm{FPR}'_g&=\mathrm{FPR}_g+\mathrm{TNR}_g\kappa_g^-,&
 \mathrm{TNR}'_g&=\mathrm{TNR}_g(1-\kappa_g^-).
\end{align*}

\begin{proof}
Conditional on $Y=1,G=g$, a final positive is either an initial positive or
an initial negative followed by appeal and reversal.  These events are
disjoint, and the latter has probability
$\mathrm{FNR}_g\alpha_g^+\rho_g^+=\mathrm{FNR}_g\kappa_g^+$.
Complementation gives the FNR identity.  Repeating the argument in the
$Y=0$ stratum gives FPR and TNR.
\end{proof}

\paragraph{Corollary 1 (positive-decision rate).}
With $\pi_g=\Prb(Y=1\mid G=g)$ and
$r_g=\Prb(D_0=1\mid G=g)$,
\[
r'_g=r_g+\pi_g\mathrm{FNR}_g\kappa_g^+
 +(1-\pi_g)\mathrm{TNR}_g\kappa_g^-.
\]

\begin{proof}
Condition on $Y$ and substitute Proposition 1.
\end{proof}

\paragraph{Proposition 2 (two-sided transformation).}
If $q_g^y=\Prb(D_1=0\mid D_0=1,Y=y,G=g)$, then
\[
 \mathrm{TPR}'_g=t_g(1-q_g^1)+(1-t_g)\kappa_g^+,
 \qquad
 \mathrm{FPR}'_g=f_g(1-q_g^0)+(1-f_g)\kappa_g^-.
\]

\begin{proof}
Partition on $D_0$ within each label stratum and apply the two rows of the
transition kernel.
\end{proof}

\subsection{Fairness preservation and signed disparity}

\paragraph{Theorem 2 (equal-opportunity preservation).}
If $t_A=t_B=t$, then post-contestation equal opportunity holds if and only if
$(1-t)(\kappa_A^+-\kappa_B^+)=0$.  When $t<1$, this is equivalent to
$\kappa_A^+=\kappa_B^+$; when $t=1$, it is automatic.

\begin{proof}
Subtract the two TPR identities in Proposition 1.
\end{proof}

\paragraph{Theorem 3 (FPR-parity preservation).}
If $f_A=f_B=f$, then post-contestation FPR parity holds if and only if
$(1-f)(\kappa_A^--\kappa_B^-)=0$.

\begin{proof}
Subtract the two FPR identities in Proposition 1.
\end{proof}

\paragraph{Corollary 2 (equalized odds).}
Under initial equalized odds, final equalized odds holds exactly when both
conditions in Theorems 2 and 3 hold.

\begin{proof}
Equalized odds is the conjunction of equal TPR and equal FPR.
\end{proof}

\paragraph{Theorem 4 (demographic-parity preservation).}
If $r_A=r_B$ initially, demographic parity is preserved exactly when
\[
\pi_A\mathrm{FNR}_A\kappa_A^+
 +(1-\pi_A)\mathrm{TNR}_A\kappa_A^-
=\pi_B\mathrm{FNR}_B\kappa_B^+
 +(1-\pi_B)\mathrm{TNR}_B\kappa_B^-.
\]

\begin{proof}
Subtract the two identities in Corollary 1 and use $r_A-r_B=0$.
\end{proof}

\paragraph{Proposition 3 (universal one-sided preservation).}
At a fixed common initial TPR $t<1$, a pair of one-sided operators preserves
equal opportunity for every admissible group distribution if and only if the
two $Y=1$ adverse-to-positive transition probabilities agree.  The analogous
claim holds for FPR at a common $f<1$.

\begin{proof}
Sufficiency follows from Theorem 2.  Necessity follows because the multiplier
$1-t$ (or $1-f$) is strictly positive.
\end{proof}

\paragraph{Theorem 5 (exact signed-gap dynamics).}
Let $d=t_A-t_B$ and
$c=(1-t_A)\kappa_A^+-(1-t_B)\kappa_B^+$.  Then $d'=d+c$ and
\begin{align*}
|d'|>|d|&\iff c(2d+c)>0,\\
|d'|<|d|&\iff c(2d+c)<0,\\
|d'|=|d|&\iff c=0\ \text{or}\ c=-2d.
\end{align*}
Direction reverses exactly when $d(d+c)<0$, and final equality holds exactly
when $c=-d$.

\begin{proof}
The identity $d'=d+c$ follows from Proposition 1.  Since absolute gaps are
nonnegative,
$|d+c|^2-|d|^2=c(2d+c)$ has the same sign as their difference.  The remaining
claims follow directly from the signs of $d,d+c$ and from $d+c=0$.
\end{proof}

\paragraph{Theorem 6 (tight attainable bounds).}
For unrestricted independent $\kappa_A^+,\kappa_B^+\in[0,1]$,
\[
d'\in[t_A-1,1-t_B].
\]
The minimum attainable absolute gap is zero when this interval contains zero
and otherwise is the distance from the interval to zero; the maximum is the
larger absolute endpoint.

\begin{proof}
Writing $a=1-t_A$ and $b=1-t_B$, the correction $a\kappa_A^+-b\kappa_B^+$
ranges continuously over $[-b,a]$.  Adding $d$ gives the displayed interval.
The absolute-value extrema follow from elementary interval geometry and are
attained at a corner or a feasible root.
\end{proof}

\paragraph{Corollary 3 (initially fair case).}
If $t_A=t_B=t$, then
$|d'|=(1-t)|\kappa_A^+-\kappa_B^+|\le1-t$, and the bound is tight.

\subsection{Endogenous access and the conditional impossibility}

\paragraph{Proposition 4 (threshold appeal response).}
If perceived value is deterministic at $v_g^y$ and the conditional cost CDF
is $F_g^y$, then $\alpha_g^y(u)=F_g^y(v_g^y+u)$.  For random value $V$,
\[
\alpha_g^y(u)=
\E\!\left[F_{C\mid V,g,y}(V+u\mid V,G=g,Y=y,D_0=0)\right].
\]

\begin{proof}
The appeal event is $\{C\le V+u\}$.  Condition on the declared stratum and,
in the random-value case, condition once more on $V$.
\end{proof}

\paragraph{Proposition 5.} \emph{Comparative statics.}
The appeal response is nondecreasing and right-continuous in assistance.  If
the conditional cost law has density $f_{C\mid V,g,y}$ and differentiation
under the expectation is valid, then
\[
(\alpha_g^y)'(u)=
\E[f_{C\mid V,g,y}(V+u\mid V,g,y)]\ge0.
\]

\begin{proof}
The indicator of $C\le V+u$ is pointwise nondecreasing in $u$, and its
expectation is a shifted CDF mixture.  Right-continuity follows from CDF
right-continuity and bounded convergence.  Under the stated regularity,
differentiation under the expectation yields the density formula.
\end{proof}

\paragraph{Theorem 7 (formal symmetry versus effective equality).}
Under common assistance $u$, deterministic perceived value $v$, and common
review success $\rho>0$,
\[
\kappa_A^+(u)-\kappa_B^+(u)
=\rho\{F_A^+(v+u)-F_B^+(v+u)\}.
\]
Thus effective correction parity holds exactly when the two conditional CDFs
agree at the operative threshold.

\begin{proof}
Apply Proposition 4 and multiply each appeal probability by the common
review-success rate.
\end{proof}

\paragraph{Theorem 8 (group-blind policy impossibility).}
Assume initial $t_A=t_B=t<1$, group-blind $u$ in a feasible set $U$, common
$v$, common $\rho>0$, and
$F_A^+(v+u)<F_B^+(v+u)$ for every $u\in U$.  No feasible group-blind policy
preserves equal opportunity.

\begin{proof}
Theorem 7 gives $\kappa_A^+(u)<\kappa_B^+(u)$.  Theorem 2 then gives
\[
t'_A-t'_B=(1-t)\rho\{F_A^+(v+u)-F_B^+(v+u)\}<0.
\]
The strictness assumptions are essential: saturation, a crossing point,
zero review success, a perfect initial classifier, unequal review quality, or
group-specific assistance can each invalidate the conclusion.
\end{proof}

\subsection{Limited-resource allocation}

Define $\mathcal{R}_g(u_g)=
\mathrm{FNR}_g[1-\rho_g^+\alpha_g^+(u_g)]$.

\paragraph{Theorem 9 (existence).}
On a nonempty compact product domain, if every $\mathcal{R}_g$ and resource cost $c_g$
is continuous and the budget set is nonempty, both the minimax residual-risk
program and the stated welfare-plus-disparity program attain optima.

\begin{proof}
The budget sublevel set is closed in a compact domain, hence compact.  Finite
maxima, sums, and absolute differences of continuous functions are continuous.
Weierstrass' theorem gives attainment.
\end{proof}

\paragraph{Theorem 10 (convexity conditions).}
If every appeal response is concave, each cost is convex, and the domain is
convex, then every $\mathcal{R}_g$ is convex and the minimax problem is convex in
epigraph form.  The separable welfare objective is convex.  A disparity-
penalized objective is convex only when its pairwise absolute residual-risk
differences are convex (for example, under affine residual risks).

\begin{proof}
$\mathcal{R}_g$ is a constant plus a nonpositive multiple of a concave response, hence
convex.  Finite maxima preserve convexity, and convex cost sublevel sets are
convex.  Weighted sums preserve convexity.  The final qualification is needed
because an absolute difference of two general convex functions need not be
convex.
\end{proof}

\paragraph{Theorem 11 (KKT marginal rule).}
For the differentiable convex separable-welfare program under Slater's
condition, there are nonnegative multipliers $\eta,\ell_g,h_g$ satisfying
\[
w_g\mathcal{R}'_g(u_g)+\eta c'_g(u_g)-\ell_g+h_g=0
\]
and complementary slackness.  If the budget binds,
$0<u_g<\bar u_g$, and $c'_g(u_g)>0$, then
\[
-\frac{w_g\mathcal{R}'_g(u_g)}{c'_g(u_g)}=\eta.
\]

\begin{proof}
Under convexity and Slater's condition, KKT conditions are necessary and
sufficient.  Interior box constraints set $\ell_g=h_g=0$; division by the
positive marginal resource cost gives the displayed equality.  Boundary
groups obey inequalities rather than this equality.
\end{proof}

\paragraph{Proposition 6 (minimum budget for a target).}
If each $\mathcal{R}_g$ is continuous and strictly decreasing over its effective range
and each $c_g$ is nondecreasing, define
$u_g^{\min}(\epsilon)=\inf\{u:\mathcal{R}_g(u)\le\epsilon\}$.  The common target is
feasible exactly when every set is nonempty and
$B\ge\sum_gc_g(u_g^{\min})$.

\begin{proof}
Continuity makes each nonempty target sublevel closed, so the infimum is
attained.  Strict monotonicity makes every feasible component at least its
minimum.  Nondecreasing costs make the sum a necessary budget; choosing the
componentwise minima proves sufficiency.
\end{proof}

For $\mathrm{FNR}_g>0$ and $\rho_g>0$, the untruncated required appeal
probability is
$p_g^*=(1-\epsilon/\mathrm{FNR}_g)/\rho_g$.  If $p_g^*>1$, equivalently
$\epsilon<\mathrm{FNR}_g(1-\rho_g)$, the target is infeasible even with
universal appeal.  This boundary must not be hidden by clipping $p_g^*$.

\paragraph{Theorem 12 (discrete allocation complexity).}
With binary-encoded nonnegative integer intervention costs and benefits, the
threshold decision version of welfare-maximizing indivisible assistance is
weakly NP-complete, even for one group.

\begin{proof}
Map an arbitrary 0--1 knapsack instance identically to individuals with the
same costs and encoded expected benefits, the same budget, and the same target
benefit.  Feasible subsets and objective values coincide.  Membership in NP
follows by checking a proposed subset.  The inherited hardness is weak, not
strong.
\end{proof}

\subsection{Identification and statistical auditing}

\paragraph{Theorem 13 (appeal-log non-identification).}
Without restrictions on appeal selection, ordinary appeal logs do not identify
the final FNR even with perfect review and observed appellant labels.

\begin{proof}
Let every case have $D_0=0$, fix
$\Prb(A=1)=a\in(0,1)$ and
$\Prb(Y=1\mid A=1)=p\in(0,1)$, and let perfect review set $D_1=Y$ for
appellants.  Non-appellants keep $D_1=0$.  Vary
$q=\Prb(Y=1\mid A=0)$ over $[0,1]$.  Since $Y$ is unobserved exactly in
the non-appellant stratum, all values of $q$ yield the same observed law, but
\[
\mathrm{FNR}'(q)=\frac{(1-a)q}{ap+(1-a)q}
\]
is strictly increasing in $q$.
\end{proof}

\paragraph{Corollary 5 (fairness non-identification).}
Applying Theorem 13 separately to two groups gives observationally equivalent
joint laws with different signed and absolute equal-opportunity gaps.  With
positive fixed appellant-positive mass, a unit gap is approached but not
attained.

\paragraph{Theorem 14 (sharp no-assumption FNR bounds).}
Let
$s=\Prb(D_0=1,Y=1\mid g)$,
$m=\Prb(D_0=0,A=1,Y=1\mid g)$,
$n=\Prb(D_0=0,A=0\mid g)$,
let $\rho$ be the identified mean reversal probability among appealed
positives, and let $z\in[0,n]$ be the latent positive mass among adverse
non-appellants.  Then
\[
F(z)=\frac{m(1-\rho)+z}{s+m+z}.
\]
When $s+m>0$ and $s+\rho m>0$, its sharp identified set is
\[
\left[\frac{m(1-\rho)}{s+m},
\frac{m(1-\rho)+n}{s+m+n}\right].
\]
If $s+m>0$ but $s+\rho m=0$, the set is $\{1\}$.  If
$s=m=0<n$, the set over compatible laws with a nonempty positive stratum is
also $\{1\}$; if $s=m=n=0$, the FNR is undefined under every compatible law.

\begin{proof}
Where defined,
\[
F'(z)=\frac{s+\rho m}{(s+m+z)^2}\ge0.
\]
Endpoint substitution gives the first interval.  The stated boundary cases
follow by direct substitution.  Sharpness holds because every feasible $z$
can be assigned inside the unobserved stratum without changing the observed
law.
\end{proof}

\paragraph{Proposition 7 (propensity-restricted sharp bounds).}
Assume $m>0$ and
$0<\underline\alpha\le\alpha\le\overline\alpha\le1$, where
$\alpha=m/(m+z)$.  Then
\[
z\in\left[
m\frac{1-\overline\alpha}{\overline\alpha},
m\frac{1-\underline\alpha}{\underline\alpha}
\right]\cap[0,n].
\]
Substitution into $F(z)$ gives sharp restricted bounds.

\begin{proof}
Solve $\alpha=m/(m+z)$ for $z=m(1-\alpha)/\alpha$.  This map is decreasing
on $(0,1]$, which reverses the two propensity endpoints.  Intersecting with
the feasible mass interval is necessary and sufficient.  If the intersection
is empty, the restriction is falsified by the observed masses.
\end{proof}

\paragraph{Theorem 15 (identification by representative audits).}
If each adverse non-appellant is audited with positive probability independently
of $Y$ conditional on the recorded stratum, and the audit reveals $Y$ without
error, then the missing prevalence $q$ and hence $z=nq$, the final FNR, and
post-contestation disparities are identified.

\begin{proof}
Conditional independence gives
\[
\Prb(Y=1\mid \text{audit}=1,D_0=0,A=0,G=g)
=\Prb(Y=1\mid D_0=0,A=0,G=g)=q_g.
\]
All remaining terms in Theorem 14 are observed.
\end{proof}

\paragraph{Proposition 8 (finite-audit contraction).}
For $M_g\ge1$ independently audited non-appellants and positive fraction
$\widehat q_g$, Hoeffding's inequality gives, with probability at least
$1-\delta$,
\[
|\widehat q_g-q_g|\le
e_g=\sqrt{\frac{\log(2/\delta)}{2M_g}}.
\]
After clipping the prevalence interval to $[0,1]$ and propagating it through
$F(n_gq)$, its width is at most
\[
\frac{2n_g(s_g+\rho_gm_g)}{(s_g+m_g)^2}e_g
\]
when $s_g+m_g>0$.

\begin{proof}
Hoeffding gives the prevalence event.  Theorem 14's map is monotone, so its
image is a confidence interval.  Its derivative with respect to $q$ is
$n_g(s_g+\rho_gm_g)/(s_g+m_g+n_gq)^2$, bounded above by the displayed
coefficient without $2e_g$.  The mean-value theorem completes the bound.
\end{proof}

\paragraph{Proposition 9 (simultaneous group bounds).}
For $K$ groups, setting
\[
e_g=\sqrt{\frac{\log(2K/\delta)}{2M_g}}
\]
gives simultaneous coverage at least $1-\delta$ for all group prevalences,
FNR images, and interval-derived pairwise gaps.

\begin{proof}
Apply Hoeffding at error probability $\delta/K$ within each group and take a
union bound.  Monotone image and interval arithmetic preserve the joint
coverage event.
\end{proof}

\subsection{Audit disposition and code correspondence}

The proof audit found no contradiction in the central signed-gap,
preservation, endogenous-access, allocation, or non-identification results.
It did require three substantive clarifications: Theorem 1 is now stated on a
signed-perturbation space rather than inferred from an unspecified probability
domain; Proposition 6 exposes infeasible targets instead of clipping the
required appeal probability; and Theorem 14/Proposition 7 now handle zero-mass
boundaries explicitly.  Regression tests cover the new identification
boundaries.  The theorem-to-code correspondence remains diagnostic rather
than a substitute for proof.
